% RESEARCH_DRAFT_SHA256: c5f2db21cb83f86f6a34fc59a3a8439f81c4a563a25a00b38f7500d385af75b2
% Layout and source-reference wrappers expanded for the web renderer.
\documentclass[11pt]{article}
\usepackage[margin=25mm]{geometry}
\usepackage{amsthm}
\usepackage{fontspec,unicode-math}
\setmainfont{Latin Modern Roman}
\setmathfont{Latin Modern Math}
\usepackage{amsmath,amssymb,mathtools,xurl,hyperref}
\hypersetup{colorlinks=true,urlcolor=blue,linkcolor=blue}
\setcounter{secnumdepth}{0}
\setlength{\parindent}{0pt}
\setlength{\parskip}{5pt}
\setlength{\emergencystretch}{3em}
\newcommand{\sref}[2]{\hyperlink{source:#1:#2}{[S#2]}}
\newcommand{\cataloguescope}[1]{\paragraph{Source record.}#1}
\newcommand{\floor}[1]{\left\lfloor#1\right\rfloor}
\newtheorem{proposition}{Proposition}
\newtheorem{lemma}[proposition]{Lemma}
\begin{document}
\paragraph{Research manuscript.} OpenAI \texttt{gpt-6-astra}, reasoning \texttt{ultra}; 11 October 2026.
\paragraph{Preserved source body.} The complete catalogue mathematical text and sources follow with layout and source links expanded, including source qualifications and its historical preparation notes. New research begins at the section headed Research attempt.

\section{Coefficient Compression and Structured Proof-System Lower Bounds}\label{statement}

\subsection{Definitions and mathematical statement}

\paragraph{Shared notation.}Standard number systems are $\mathbb N,\mathbb Z,\mathbb Q,\mathbb R,\mathbb C$. Specialized notation and index ranges follow the cited source and the definitions below.


A polynomial is a sum of squares (SOS) if it is $\sum_j q_j^2$ for polynomials $q_j$. Moment-SOS hierarchies impose positivity through increasingly large truncated moment matrices and corresponding degree-bounded SOS certificates.

A refutation is a derivation of contradiction from an unsatisfiable constraint system using specified inference rules. Proof size counts the encoded derivation, and degree in algebraic proof systems is the largest polynomial degree appearing.
\paragraph{Record scope and qualifications.}
The catalogue title is retained as metadata. Named certificates, contracts, bounds or frameworks in that title are not asserted as results of the original source extract.
Source-record limitation: The combined target omits the exact coefficient-compression preservation criterion and the inference rules of the proposed DNNF proof system; importing generic definitions would not recover the mathematical problem.
\paragraph{Statement recorded in the supplied source.}
Can the coefficient bit length of every $n$-variable linear threshold function be reduced efficiently from $n^{100}$ bits to $n^2$ bits? Can one prove Polynomial Calculus, Sherali–Adams, or Sum-of-Squares degree lower bounds for graph-isomorphism formulas $\operatorname{ISO}(G,H)$ beyond CFI-type graphs? How strong is a deductive proof system whose lines are DNNF formulas? \href{https://doi.org/10.4171/owr/2024/15}{[S2]}
\paragraph{Source status.}
The supplied catalogue labels this source record open; that metadata was checked against the supplied snapshot on 2026-10-11. The label is preserved for provenance and is not an independent verification of present open status or mathematical correctness. The catalogue literature assessment is dated 2026-08-22.
\subsection{Short English statement}
The source asks about coefficient Compression and Structured Proof-System Lower Bounds.
\subsection{Sources}
\begin{itemize}
\item[S1] Ulam AI and the UnsolvedMath Contributors, UnsolvedMath: A Curated Collection of Open Mathematics Problems, catalogue entry. Catalogue curation and metadata: CC BY 4.0. The source extract is retained; definitions, English summaries and qualifications are editorial additions.
\url{https://huggingface.co/datasets/ulamai/UnsolvedMath}.
\item[S2] Proof Complexity and Beyond (2024), Oberwolfach Reports, DOI 10.4171/owr/2024/15, pp. 871--934
\url{https://doi.org/10.4171/owr/2024/15}.
\item[S3] Oberwolfach Report 15/2024
\url{https://oa.tib.eu/renate/bitstreams/b0ef46dc-1bdf-45b7-a4f5-ca0a9e60a7d6/download}.
\item[S4] Micun and Berkholz, Proof Systems Based on Structured Circuits
\url{https://arxiv.org/abs/2605.12378}.
\item Source-specific Oberwolfach Reports licences, including CC BY-SA 4.0 where applicable, are retained; see THIRD\_PARTY\_NOTICES.md.
\end{itemize}
\paragraph{Preparation and review.}
The original source and UnsolvedMath contributors retain attribution for the source extract. Codex (GPT-6; exact serving revision not exposed) assisted with definitions, English summaries, scope notes and formatting. No new human mathematical authorship or independently verified solution is asserted. Catalogue identity checked against the complete supplied registry and current repository. Distinct upstream records, including repeated questions and broad or damaged extracts, are retained for complete coverage of the source CS category. A repeated record is not a new independent conjecture. Complete mathematical-equivalence and current-literature review has not been performed. Preparation and final review follow contribution prompts C and D, with source gaps disclosed. Structural and rendering checks do not establish a new theorem or open status.
\label{end:statement}

\section{Research attempt}
\subsection{Literature, source context and scope}
The full original catalogue body and source qualifications are preserved above. The relevant primary workshop passages were read in the cached publisher report. Historical claims and subsequent attributed results are distinguished from the mathematical work below. A restricted example, standard lemma or known theorem is not counted as new progress on the entire original target. Missing assumptions and repeated or rhetorical extracts are not silently turned into different solved problems.

\subsection{Existence of small threshold weights does not compress them efficiently}
The source itself notes the existence of small threshold representations and asks for an efficient reduction from very long given coefficients. We spell out that existence argument to locate its computational gap. Consider a linear threshold function on $\{0,1\}^n$, $n\ge1$, with ties labeled positive. A representation is a vector $w$ and threshold $t$ deciding whether $w\cdot x-t\ge0$.
\begin{proposition}
Every such function has an integer representation whose $n+1$ coefficients each have absolute value at most $(n+1)^{(n+1)/2}$. This is $O(n\log(n+1))$ bits per coefficient, or $O(n^2\log(n+1))$ total bits.
\end{proposition}
\begin{proof}
If there are negative inputs, decrease the threshold by less than their minimum positive distance from the threshold. Every positive input then has positive margin and every negative input stays negative. If there are no negative inputs, choose the constant-positive representation. Scaling gives signed margins at least one. For labels $y_x\in\{-1,1\}$, the polyhedron in $d=n+1$ variables $z=(w,t)$ defined by
\[
y_x(x,-1)\cdot z\ge1\qquad(x\in\{0,1\}^n)
\]
is therefore nonempty. Its constraint normals span $\mathbb R^d$, since multiplying by nonzero signs preserves the span of $(x,-1)$, which contains a basis from $x=0,e_1,\ldots,e_n$.
It has a point with $d$ independent tight constraints. To see this directly, start at a feasible point. If tight normals do not span the space, move along a nonzero vector orthogonal to all of them. Some remaining normal has nonzero inner product with this direction. Choose a sign in which at least one slack decreases; the first tight constraint is reached after a finite positive step, preserving feasibility. Its normal is independent of the old tight normals, since it is not orthogonal to the direction. Repeat at most $d$ times.
At the resulting point, solve the independent tight equations $Az=\mathbf1$. Entries of $A$ lie in $\{-1,0,1\}$. Cramer's rule gives $z_j=\det A_j/\det A$, where $A_j$ replaces column $j$ by ones. Each determinant has absolute value at most $d^{d/2}$: the determinant is the product of lengths of successive orthogonal components of its rows, each at most the row norm $\sqrt d$. Multiply $z$ by $|\det A|$, a positive integer. All coordinates become integers of the stated size and all margins remain positive. Thus the labels are preserved.
\end{proof}
The construction potentially uses all $2^n$ labeled inputs to identify tight constraints. Enumerating them or their $d$-tuples is not polynomial-time compression from the succinct original weights. The existence proof therefore stops before the requested algorithm. The source's per-coefficient bit convention must also not be replaced by a total-bit claim.
Naive coarse rounding is insufficient: weights $(M,M+1)$ and threshold $M+1$, with $M\ge1$, label $(1,0)$ negative and $(0,1)$ positive. Replacing both weights by the same rounded value makes these two inputs indistinguishable under every threshold.
\subsection{Verification and remaining gaps}
For all threshold-representable labelings in two variables, integer active-constraint solutions were enumerated and checked against all four inputs; the determinant and rounding bounds were verified exactly. No efficient search rule follows. The graph-isomorphism degree question and the DNNF proof-system question are independent components requiring their precise systems; they are not settled by threshold geometry. The cited 2026 structured-circuit paper's primary landing page was consulted for context, not as a proof of the whole bundled target.

\section{References and source attribution}
Primary workshop report: \url{https://doi.org/10.4171/owr/2024/15}. Relevant passages read: 2024/15 p.876, threshold-coefficient compression, non-CFI isomorphism degree and DNNF proof-system questions.

\url{https://arxiv.org/abs/2605.12378}. Opened the primary structured-circuit paper landing page for context; no full proof or broad DNNF-system equivalence is claimed verified.

All retrievals were checked on 11 October 2026. Original question attribution remains in the preserved Sources section. External theorems are attributed and not claimed as new discoveries or fully verified merely by citation.


\paragraph{Integration review.} Codex (GPT-6; exact serving revision not exposed) reviewed the written arguments and expanded layout and source-reference wrappers. Repository metadata and the production web renderer were checked using Python 3.13 and JavaScript with MathJax 3.2.2. This model-assisted review is not independent human proof verification; no PDF compilation or formal proof checking is asserted.

\section{Completion Estimate}
\noindent\textbf{COMPLETION: 0\%}
No substantiated new progress toward the entire original target is claimed. The worked results and failed route are delimited above; a known result, illustrative model, or auxiliary obstruction is not counted as original-target completion.

\section{Final status and attribution}
\textbf{Attempt status: unresolved.} This describes this research attempt, not a new ruling on the source catalogue's status or on external literature claims.
The mathematics and this self-audit were generated by OpenAI \texttt{gpt-6-astra}, reasoning \texttt{ultra}, on 11 October 2026. Codex (GPT-6; exact serving revision not exposed) orchestrated the batch, read and reviewed its written mathematical arguments, and contributed formatting/repository review. No human mathematical authorship, independent human proof checking, or formal verification is claimed. The historical preparation paragraphs in the reproduced catalogue body concern that body only.

Contribution prompt A was applied to the complete definition and supplied source record, followed by the adversarial checks of prompt B. Local reproduction files are \texttt{batch025/checks.py}, \texttt{computations.json}, and the per-ID audit JSON. They were not public links at preparation time. Numerical checks are finite and do not prove unrestricted results. No LaTeX engine was found locally; compilation was not performed by this research agent.

\end{document}
